\documentclass[12pt,a4paper]{article}
\usepackage[utf8]{inputenc}
\usepackage{amsmath}
\usepackage{amsfonts}
\usepackage{amssymb}
\usepackage{geometry}
\usepackage{booktabs}
\usepackage{cite}
\usepackage{microtype} % Handles spacing to reduce Overfull \hbox warnings

\geometry{margin=1in}

\begin{document}

\begin{center}
    \large \textbf{The Mechanical Vacuum: \\ Frame-Dragging, Satiation Velocity, and the Variable Gravitational Constant} \\
    \vspace{1em}
    \normalsize Robin Skavberg \\
    \normalsize February 2026
\end{center}

\begin{abstract}
Standard General Relativity (GR) treats $G$ as a static scalar constant. This paper proposes an \textbf{Equilibrium Perspective}, wherein $G$ represents the dynamic restorative response of the vacuum field. We derive $G$ as a function of vacuum restorative stiffness and matter-energy density $\rho_m$. By defining a state of \textbf{Mechanical Satiation}, we provide a physical basis for the Newtonian limit and the behavior of high-density singularities, while offering a mechanical framework that complements and expands upon the dark matter paradigm in low-density galactic halos. 
\end{abstract}

\section{The Mechanical Basis of Gravitation}
In this framework, the gravitational constant $G$ is an emergent property of the vacuum's coupling efficiency. We define $c^2$ as the \textbf{intrinsic restorative stiffness} of the vacuum medium, while the matter-energy density $\rho_m$ acts as a mechanical load.

The foundational expression for the gravitational constant $G$ is derived as:
\begin{equation}
G = \frac{c^2}{8\pi\rho_m}
\end{equation}
Here, $\rho_m$ is interpreted as an effective linear mass density ($M/L$). Dimensionally, this ensures that $G$ maintains the units $\text{m}^3 \cdot \text{kg}^{-1} \cdot \text{s}^{-2}$. This relationship suggests that as matter populates a region, the vacuum's \textbf{tensor susceptibility} decreases—the field becomes ``stiffer'' and less responsive to additional displacement.

% Added manual break to fix Overfull hbox at line 34
\section{Tensor Saturation and the \\ Critical Density Threshold}
To maintain consistency with Solar System observations and extreme astrophysical phenomena, the model must account for the observed invariance of $G$ in high-density regimes and its behavior in the presence of compact objects.

\subsection{The Asymptotic Satiation Formula}
The transition from the high-susceptibility ground state ($G_{vac}$) to the satiated Newtonian limit ($G_N$) is governed by a critical density threshold $\rho_{crit}$:
\begin{equation}
G(\rho_m) = G_{vac} \cdot \exp\left(-\frac{\rho_{local}}{\rho_{crit}}\right) + G_{N}
\end{equation}
In this regime, $\rho_{crit}$ represents the point at which the vacuum's tensor flexibility is exhausted.

\subsection{The Milgromian Link ($a_0$)}
A rigid mathematical link is established by correlating $\rho_{crit}$ with the Milgromian acceleration constant, $a_0 \approx 1.2 \times 10^{-10} \text{ m/s}^2$. We define the critical density as the state where the vacuum's restorative pressure is in equilibrium with the field energy density at acceleration $a_0$:
\begin{equation}
\rho_{crit} = \frac{a_0^2}{8\pi G_N^2 \sigma}
\end{equation}
where $\sigma$ represents the surface tension of the vacuum field.

\subsection{Absolute Satiation: The Black Hole Limit}
In the case of a Black Hole, the matter-energy density $\rho_m$ exceeds the critical threshold $\rho_{crit}$ by several orders of magnitude. In this mechanical model, the Schwarzschild radius $R_s = 2GM/c^2$ represents the boundary of \textbf{Absolute Mechanical Satiation}. 

As $\rho_m \to \infty$ at the horizon, the vacuum impedance $Z_v$ approaches an infinite limit, and the coupling $G$ is clamped to its absolute minimum. This suggests that the ``singularity'' is not a point of infinite curvature, but a region of \textbf{Mechanical Stagnation} where the restorative stiffness of the vacuum is entirely overwhelmed by the load, effectively freezing the propagation of waves ($v_\gamma \to 0$) within the satiated volume.

\subsection{The Theoretical Ceiling: Ground-State Maximum ($G_{vac}$)}
The gravitational coupling constant $G$ does not diverge to infinity as $\rho_m \to 0$. Instead, it is bounded by the intrinsic energy density of the vacuum ground state, $\rho_{vac}$. We define the maximum coupling strength as:
\begin{equation}
G_{vac} = \frac{c^2}{8\pi \rho_{vac}}
\end{equation}
In this ground state, the vacuum exhibits maximum tensor susceptibility. The transition from the Newtonian floor ($G_N$) in satiated regions to the vacuum ceiling ($G_{vac}$) in cosmic voids accounts for the observed anomalies in large-scale structure and expansion rates without the necessity of introducing dark energy as a separate fluid.

\section{Derivation of the Propagation Floor ($v_{\gamma}$)}
The link between the gravitational coupling ($G$) and the phase velocity of waves in the field ($v_{\gamma}$) defines the dynamic range of the vacuum.

\subsection{Vacuum Impedance and Refractive Index}
We model the vacuum's \textbf{mechanical impedance} $Z_v$ as a function of the matter-energy load. The effective refractive index $n_v$ is determined by:

\begin{equation}
n_v = \sqrt{1 + \rho_m} = \sqrt{1 + \frac{c^2}{8\pi G}}
\end{equation}

\subsection{Defining the Satiation Ratio ($\mathcal{S}$)}
The maximum wave speed in an unsatiated ground state is $c$. As density increases toward satiation, the phase velocity drops toward a \textbf{Propagation Floor} ($v_{\gamma}$), which is the minimum wave speed permissible under maximum mechanical load:
\begin{equation}
v_{\gamma} = \frac{c}{n_v} = c \sqrt{\frac{8\pi G}{8\pi G + c^2}}
\end{equation}
Using SI values for $G_N$ and $c$, we find the satiation ratio $\mathcal{S} = v_{\gamma}/c \approx 1.3 \times 10^{-13}$, yielding a local propagation floor of approximately $40 \, \mu\text{m/s}$.

\section{Correspondence with Screening Mechanisms}
This mechanical satiation provides a physical foundation for established screening theories.

\textbf{Chameleon Screening:} Interpreted as the non-linear increase in the vacuum's effective elastic modulus in high-density regions, rendering the field unresponsive to further displacement.

\textbf{Symmetron Theory:} The spontaneous symmetry breaking at low densities is viewed as the transition of the vacuum from a satiated, rigid state to a high-susceptibility phase where the vacuum expectation value (VEV) shifts.

\textbf{DGP Vainshtein Screening:} The 4D/5D transition is recontextualized as the ``satiation bottleneck.'' Near massive sources, the 4D manifold becomes so stiff that it restricts higher-dimensional flux leakage, forcing Newtonian behavior.

% Manual break to fix Overfull hbox at line 88
\section{Evidentiary Observations and \\ Empirical Alignment}
The mechanical vacuum model and the derived propagation floor $v_{\gamma}$ align with several key astrophysical observations:

\subsection{The Hubble Tension: Dynamic Expansion}
The discrepancy in $H_0$ results from the vacuum's shifting restorative stiffness as $\rho_m$ evolves. Local measurements reflect an unsatiated regime where $G > G_N$:
\begin{equation}
H(z) = H_{0} \sqrt{\Omega_m (1+z)^3 \left( \frac{G(z)}{G_N} \right) + \Omega_\Lambda}
\end{equation}

% Using \sloppy here to fix the spacing issues in this specific paragraph
{\sloppy
\textbf{Mathematical Proof with Observed Values:}
Standard observations yield $H_{local} \approx 73.2$ km/s/Mpc and $H_{CMB} \approx 67.4$ km/s/Mpc. In this framework, the expansion rate scales with the vacuum coupling ratio:
The correlation between the $8.6\%$ Hubble Tension and the $8.6\%$ vacuum recovery is derived from the power-law relationship between the gravitational coupling $G$ and the expansion rate of the universe $H$. \par}

\subsubsection*{The Relationship Between $H$ and $G$}
In standard Friedmann cosmology, the Hubble parameter $H$ scales with the square root of the gravitational constant:
\[ H \propto \sqrt{G} \]

\subsection{Deriving the $8.6\%$ Expansion Differential}

\section*{1. Increase from baseline}
\[ \frac{73 - 67}{67} \times 100 \approx \mathbf{8.96\%} \]

\section*{2. Difference relative to observed value}
\[ \frac{73 - 67}{73} \times 100 \approx \mathbf{8.2\%} \]

\section*{3. Symmetric (midpoint) percentage difference}
\[ \frac{2(73 - 67)}{73 + 67} \times 100 \approx \mathbf{8.6\%} \]

The $8.6\%$ represents the midpoint ``Tension'' observed between: 
\textbf{The Baseline ($H_{CMB}$):} $67.4$ km/s/Mpc and 
\textbf{The Measured ($H_{local}$):} $73.2$ km/s/Mpc.

\subsection{Local Frame-Dragging (Gravity Probe B)}
The derived propagation floor matches the magnitude of linear frame-dragging measured locally.
\begin{equation}
v_{drift} = \Omega_{LT} \cdot R_{\oplus} \approx 4.2 \times 10^{-5} \text{ m/s} \approx v_{\gamma}
\end{equation}

\section{Conclusion}
The Mechanical Vacuum model provides a refined framework that reinterprets $G$ as a dynamic equilibrium state. This shift bridges the gap between local precision measurements and galactic-scale observations.

\section*{Author Information}
\noindent \textbf{AI Statement:} LLM technology was utilized as a technical collaborator for LaTeX typesetting and document structure optimization. Unique physical theories were produced by the human author.

\begin{thebibliography}{9}
\bibitem{1} Einstein, A. (1916). Annalen der Physik, 354(7), 769--822.
\bibitem{2} Skavberg, R. (2026). The Mechanical Resolution of the Gravitational Constant. Zenodo.
\bibitem{3} Everitt, C. W. F., et al. (2011). Gravity Probe B. Phys. Rev. Lett., 106(22), 221101.
\bibitem{4} Unruh, W. G. (1976). Phys. Rev. D, 14(4), 870.
\bibitem{5} Sanchez, C., et al. (2017). Cosmic voids and lensing anomalies. JCAP, 2017(09), 015.
\end{thebibliography}

\end{document}